% from ega4-12-fr.tex
\oldpage[IV-3]{177}
  (ii)) où l'on prend pour $Z'$ une des deux composantes irréductibles de
  $X=\operatorname{Spec}(B)$. Il est immédiat que ce morphisme restriction
  n'est pas ouvert aux points de $Z'$ au-dessus du «~point double~» de $Y'=Y$.
  \item Dans les hypothèses de (12.1.1.5), avec
  $\mathcal F=\mathcal O_X$, on ne peut affaiblir la condition «~$f$ est
  plat~» en «~$f$ est universellement ouvert~» (cf. (2.4.6)), comme nous le
  verrons plus loin sur un exemple (14.4.10, (i)).
\end{enumerate}
\end{remarks}

\begin{corollary}[12.1.3]
\label{IV.12.1.3-fr}
Sous les hypothèses de (12.1.1), la fonction
$x\mapsto\operatorname{coprof}((\mathcal F_{f(x)})_x)$ est semi-continue
supérieurement dans $X$ et la fonction
$x\mapsto\operatorname{prof}_x^*(\mathcal F_{f(x)})$ (10.8.1) est
semi-continue inférieurement dans $X$.
\end{corollary}

\begin{proof}
La première assertion n'est autre qu'une formulation équivalente de (12.1.1,
(v)). Pour prouver la seconde, on peut d'abord, compte tenu de (10.8.7) et
(10.8.8), se ramener comme dans (12.1.1) au cas où $Y$ est le spectre d'un
anneau de valuation discrète $A$ de point fermé $y$ et de point générique $y'$,
$X=\operatorname{Spec}(B)$ étant affine, $\mathcal F=\widetilde M$, où $M$ est
un $B$-module de type fini. Puisque $\mathcal F$ est $f$-plat, toute composante
irréductible de $\operatorname{Supp}(\mathcal F)$ qui contient $x\in X_y$
domine $Y$ (2.3.4), autrement dit son point générique $z'$ appartient à
$X_{y'}$~; les composantes irréductibles $Z$ de
$\operatorname{Supp}(\mathcal F)$ qui contiennent une générisation $x'$ de
$x$ appartenant à $X_{y'}$ sont donc exactement celles qui contiennent $x$ et
en outre, d'après (12.1.1.5), on a $\dim(Z_{y'})=\dim(Z_y)$, d'où, si
$T=\operatorname{Supp}(\mathcal F)$,
$\dim_x(T_y)=\dim_{x'}(T_{y'})$. De plus, pour une uniformisante $t$ de $A$, on
a vu dans (12.1.1, (v)) que l'on a
$\operatorname{coprof}(M_x)=\operatorname{coprof}(M_x/tM_x)$, et puisque
d'autre part $\operatorname{coprof}(M_{x'})\leq\operatorname{coprof}(M_x)$
(6.11.5), on voit que l'on a
$\operatorname{coprof}((\mathcal F_{y'})_{x'})\leq
\operatorname{coprof}((\mathcal F_y)_x)$, la relation
$\operatorname{prof}_x^*(\mathcal F_y)\leq
\operatorname{prof}_{x'}^*(\mathcal F_{y'})$ découle alors de (10.8.7).
\end{proof}

\begin{remark}[12.1.4]
\label{IV.12.1.4-fr}
On déduit aussi de (12.1.1, (i)) que la fonction
$x\mapsto\dim_x(X_{f(x)})$ est semi-continue supérieurement dans $X$, mais
nous verrons plus loin (13.1.3) que cette propriété est vraie même sans
supposer $f$ plat.
\end{remark}

\begin{corollary}[12.1.5]
\label{IV.12.1.5-fr}
Soient $Y$ un préschéma localement noethérien, $f:X\to Y$ un morphisme
localement de type fini, $\mathcal F$ un $\mathcal O_X$-Module cohérent et
$f$-plat, $\mathcal G$ un $\mathcal O_Y$-Module cohérent, $x$ un point de $X$,
$y=f(x)$. Supposons que $\mathcal G$ ait la propriété $(S_k)$ au point $y$, et
que $\mathcal F_y$ ait la propriété $(S_k)$ au point $x$ et soit
équidimensionnel au point $x$. Alors~:
\begin{enumerate}[label=(\roman*)]
  \item $\mathcal F\otimes_{\mathcal O_Y}\mathcal G$ possède la propriété
  $(S_k)$ au point $x$.
  \item Si de plus $Y$ est localement immersible dans un schéma régulier, ou
  est un préschéma excellent (7.8.5), il existe un voisinage de $x$ dans $X$
  tel que $\mathcal F\otimes_{\mathcal O_Y}\mathcal G$ ait la propriété
  $(S_k)$ dans ce voisinage.
\end{enumerate}
\end{corollary}

\begin{proof}
En effet, par (12.1.1, (iv)), il existe un voisinage ouvert $V$ de $x$ tel que
pour tout $x'\in V$, $\mathcal F_{f(x')}$ vérifie $(S_k)$ au point $x'$.
D'autre part, $\mathcal G$ vérifie $(S_k)$ en tous les points de
$Y'=\operatorname{Spec}(\mathcal O_y)$ (5.7.2). Remplaçant $Y$ par $Y'$ et $X$
par $V\times_YY'$, on est ramené (compte tenu de ($\mathbf I$, 3.6.5)) au cas
où $\mathcal G$ vérifie $(S_k)$ dans $Y$ tout entier et où, pour tout
$y\in f(X)$, $\mathcal F_y$ possède la propriété $(S_k)$ en tous les points de
$f^{-1}(y)$~; comme $\mathcal F$ est $f$-plat, on sait alors (6.4.1) que
$\mathcal F\otimes_{\mathcal O_Y}\mathcal G$ possède la propriété $(S_k)$ dans
$X$, ce qui prouve (i). Pour démontrer (ii), il suffit d'observer que, sous les
hypothèses faites, $\mathcal G$ possède la
